This work analyzes the detection and behavior of intrusions in computer networks through the implementation of the T-Pot platform in a controlled experimental environment. For this purpose, a laboratory based on an isolated local area network (LAN) was designed, allowing the simulation of attacks without compromising real infrastructures. Using a passive observation methodology, multiple honeypots were deployed, enabling the capture and analysis of attack vectors such as port scans, brute-force attempts, and denial-of-service attacks. This facilitated the identification of behavioral patterns and attacker profiles, while demonstrating the effectiveness of T-Pot in recording and correlating events in near real time, thereby consolidating it as a useful tool for forensic analysis and the strengthening of proactive cybersecurity strategies.